\NormalFont\ShortTitle{Cantares}
{\MT Cantares de Salomão

\par }\ChapOne{1}{\PP \VerseOne{1}Cântico dos cânticos, que é de Salomão.
\VS{2}{\ADD{Ela}} : Beije-me ele com os beijos de sua boca, porque teu amor é melhor do que o vinho.
\VS{3}O cheiro dos teus perfumes é agradável; teu nome é
{\ADD{como}} perfume sendo derramado, por isso as virgens te amam.
\VS{4}Toma-me contigo, e corramos; traga-me o rei aos seus quartos.{\ADD{Moças}} : Em ti nos alegraremos e nos encheremos de alegria; nos agradaremos mais de teu amor do que do vinho;{\ADD{Ela}} : Elas estão certas em te amar;
\VS{5}Eu sou morena, porém bela, ó filhas de Jerusalém:
{\ADD{morena}} como as tendas de Quedar,
{\ADD{bela}} como as cortinas de Salomão.
\VS{6}Não fiquem me olhando por eu ser morena, pois o sol brilhou sobre mim; os filhos de minha mãe se irritaram contra mim,
{\ADD{e}} me puseram para cuidar de vinhas;
{\ADD{porém}} minha própria vinha, que me pertence, não cuidei.
\VS{7}Dize-me, amado de minha alma: onde apascentas
{\ADD{o teu gado}} ? Onde
{\ADD{o}} recolhes ao meio- dia? Para que ficaria eu como que coberta com um véu por entre os gados de teus colegas?
\FTNT{*}{{\FR 1:7 }
coberta por um véu
obscuro - trad. alt. andando sem rumo
}
\VS{8}{\ADD{Ele}} : Se tu, a mais bela entre as mulheres, não sabes, sai pelos rastros das ovelhas, e apascenta tuas cabras junto às tendas dos pastores.
\VS{9}Eu te comparo, querida, às éguas das carruagens de Faraó.
\VS{10}Agradáveis são tuas laterais da face entre os enfeites, teu pescoço entre os colares.
\VS{11}Enfeites de ouro faremos para ti, com detalhes de prata.
\VS{12}{\ADD{Ela}} : Enquanto o rei está sentado à sua mesa, meu nardo dá a sua fragrância.
\VS{13}Meu amado é para mim
{\ADD{como}} um saquinho de mirra que passa a noite entre meus seios;
\VS{14}Meu amado é para mim
{\ADD{como}} um ramalhete de hena nas vinhas de Engedi.
\VS{15}{\ADD{Ele}} : Como tu és bela, minha querida! Como tu és bela! Teus olhos são
{\ADD{como}} pombas.
\VS{16}{\ADD{Ela}} : Como tu és belo, meu amado! Como
{\ADD{tu és}} agradável! E o nosso leito se enche de folhagens.
\VS{17}As vigas de nossa casa são os cedros, e nossos caibros os ciprestes.

\par }\Chap{2}{\PP \VerseOne{1}Eu sou a rosa de Sarom, o lírio dos vales.
\VS{2}{\ADD{Ele}} : Como um lírio entre os espinhos, assim é minha querida entre as moças.
\VS{3}{\ADD{Ela}} : Como a macieira entre as árvores do bosque, assim é o meu amado entre os rapazes; debaixo de sua sombra desejo muito sentar, e doce é o seu fruto ao meu paladar.
\VS{4}Ele me leva à casa do banquete, e sua bandeira sobre mim é o amor.
\VS{5}Sustentai-me com passas, fortalecei-me com maçãs; porque estou fraca de amor.
\VS{6}{\ADD{Esteja}} sua mão esquerda abaixo de minha cabeça, e sua direita me abraça.
\VS{7}Eu vos ordeno, filhas de Jerusalém: jurai pelas corças e pelas cervas do campo, que não acordeis, nem desperteis ao amor, até que ele queira.
\VS{8}{\ADD{Esta é}} a voz do meu amado: vede-o vindo, saltando sobre os montes, pulando sobre os morros.
\VS{9}Meu amado é semelhante ao corço, ou ao filhote de cervos; eis que está atrás de nossa parede, olhando pelas janelas, observando pelas grades.
\VS{10}Meu amado me responde, e me diz:{\ADD{Ele}} : Levanta-te, querida minha, minha bela, e vem.
\VS{11}Porque eis que o inverno já passou; a chuva se acabou, e foi embora.
\VS{12}As flores aparecem na terra, o tempo da cantoria chegou; e ouve-se a voz da rolinha em nossa terra.
\VS{13}A figueira está produzindo seus figos verdes, e as vides florescentes dão cheiro; levanta-te, querida minha, minha bela, e vem.
\VS{14}Pomba minha, que andas pelas fendas das rochas no oculto das ladeiras, mostra-me tua face, faze-me ouvir tua voz; porque tua voz é doce, e tua face agradável.
\VS{15}Tomai-nos as raposas, as raposinhas, que danificam as vinhas, porque nossas vinhas estão florescendo.
\VS{16}Meu amado é meu, e eu sou sua; ele apascenta entre os lírios.
\VS{17}Antes do dia se romper, e das sombras fugirem, volta, amado meu, faze-te semelhante ao corço, ao filhote de cervos, sobre os montes de Beter.

\par }\Chap{3}{\PP \VerseOne{1}{\ADD{Ela}} :Durante as noites busquei em minha cama a quem minha alma ama; busquei-o, mas não o achei.
\VS{2}Por isso me levantarei, e percorrerei a cidade, pelas ruas e pelas praças; buscarei a quem minha alma ama; busquei-o, mas não o achei.
\VS{3}Os guardas que rondavam pela cidade me encontraram.
{\ADD{Eu lhes perguntei}} : Vistes a quem minha alma ama?
\VS{4}Pouco
{\ADD{depois}} de me afastar deles, logo achei a quem minha alma ama. Eu o segurei, e não o deixei ir embora, até eu ter lhe trazido à casa de minha mãe, ao cômodo daquela que me gerou.
\VS{5}Eu vos ordeno, filhas de Jerusalém: jurai pelas corças e pelas cervas do campo que não acordeis nem desperteis ao amor, até que ele queira.
\VS{6}Quem é esta, que sobe do deserto como colunas de fumaça, perfumada com mirra, incenso, e com todo tipo de pó aromático de mercador?
\VS{7}Eis a cama móvel de Salomão! Sessenta guerreiros estão ao redor dela, dentre os guerreiros de Israel;
\VS{8}Todos eles portando espadas, habilidosos na guerra; cada um com sua espada à cintura, para
{\ADD{o caso de haver}} um ataque repentino de noite.
\VS{9}O rei Salomão fez para si uma liteira de madeira do Líbano.
\VS{10}Suas colunas ele fez de prata, seu assoalho de ouro, seu assento de púrpura; por dentro coberta com
{\ADD{a obra do}} amor das filhas de Jerusalém.
\VS{11}Saí, ó filhas de Sião, e contemplai ao rei Salomão, com a coroa a qual sua mãe o coroou, no dia de seu casamento, no dia da alegria do seu coração.

\par }\Chap{4}{\PP \VerseOne{1}{\ADD{Ele}} : Como tu és bela, minha querida! Como tu és bela! Teus olhos por trás do véu são
{\ADD{como}} pombas; teu cabelo é como um rebanho de cabras, que descem do monte Gileade.
\VS{2}Teus dentes são como
{\ADD{ovelhas}} tosquiadas, que sobem do lavatório; todas elas têm gêmeos, e nenhuma delas é estéril.
\VS{3}Teus lábios são como uma fita de escarlate, e tua boca é bonita; tuas têmporas são como pedaços de romã por detrás do véu.
\VS{4}Teu pescoço é como a torre de Davi, construída como fortaleza; mil escudos estão nela pendurados, todos escudos de guerreiros.
\VS{5}Teus dois seios são como dois filhos gêmeos da corça, que pastam entre os lírios.
\VS{6}Antes do dia nascer, e das sombras fugirem, irei ao monte de mirra, e ao morro do incenso.
\VS{7}Tu és bela, minha querida! Não há defeito algum em ti.
\VS{8}{\ADD{Vem}} comigo do Líbano, ó esposa minha, vem comigo do Líbano; desce do cume de Amaná, do cume de Senir e de Hermom, das montanhas das leoas, dos montes dos leopardos.
\VS{9}Tomaste meu coração, minha irmã, minha esposa; tomaste o meu coração com um de teus olhos, com um colar de teu pescoço.
\VS{10}Como são agradáveis os teus amores, minha irmã, minha esposa! São bem melhores do que o vinho; e o cheiro de teus unguentos
{\ADD{é melhor}} que todas as especiarias.
\VS{11}Favos de mel descem de teus lábios, ó esposa; mel e leite estão debaixo de tua língua; e o cheiro de teus vestidos é como o cheiro do Líbano.
\VS{12}Jardim fechado és tu, minha irmã, minha esposa; manancial fechado, uma fonte selada.
\VS{13}Tuas plantas são uma horta de romãs, com frutos excelentes: hena e nardo.
\VS{14}Nardo, açafrão, cálamo e canela; com todo tipo de árvores de incenso: mirra, aloés, com todas as melhores especiarias.
\VS{15}{\ADD{Tu és}} uma fonte de jardins, um poço de águas vivas que corre do Líbano.
\VS{16}{\ADD{Ela}} : Levanta-te, vento norte! E vem, ó vento sul! Assopra em meu jardim, para que espalhem os seus aromas! Que meu amado venha a seu jardim, e coma de seus excelentes frutos.

\par }\Chap{5}{\PP \VerseOne{1}{\ADD{Ele}} : Já cheguei ao meu jardim, minha irmã,
{\ADD{minha}} esposa; colhi minha mirra com minha especiaria; comi meu mel, bebi meu vinho com meu leite.
{\ADD{Autor do poema (ou os amigos dos noivos)}} : Comei, amigos; bebei, ó amados, e embebedai-vos.
\VS{2}{\ADD{Ela}} : Eu estava dormindo, mas meu coração vigiava; era a voz de meu amado, batendo: Abre-me, minha irmã, minha querida, minha pomba, minha perfeita! Porque minha cabeça está cheia de orvalho, meus cabelos das gotas da noite.
\VS{3}Já tirei minha roupa; como
{\ADD{voltarei}} a vesti-la? Já lavei meus pés; como vou sujá-los
{\ADD{de novo}} ?
\VS{4}Meu amado meteu sua mão pelo buraco
{\ADD{da porta}} ,e meu interior se estremeceu por ele.
\VS{5}Eu me levantei para abrir a meu amado; minhas mãos gotejavam mirra, e meus dedos
{\ADD{derramaram}} mirra sobre os puxadores da fechadura.
\VS{6}Eu abri a porta ao meu amado, porém meu amado já tinha saído e ido embora; minha alma como que saía
{\ADD{de si}} por causa de seu falar.
\FTNT{*}{{\FR 5:6 }
seu falar
trad. alt. sua partida
} Eu o busquei, mas não o achei; eu o chamei, mas ele não me respondeu.
\VS{7}Os guardas que rondavam pela cidade me encontraram; eles me bateram
{\ADD{e}} me feriram; os guardas dos muros tiraram de mim o meu véu.
\VS{8}Eu vos ordeno, filhas de Jerusalém: jurai que, se achardes o meu amado, dizei a ele que estou doente de amor.
\VS{9}{\ADD{Moças}} : Ó, tu, mais bela entre as mulheres, como é o teu amado, a quem tu amas mais do que qualquer outro? Como é este teu amado, mais amado que qualquer outro, por quem assim nos mandas jurar?
\VS{10}{\ADD{Ela}} : Meu amado é radiante e rosado; ele se destaca entre dez mil.
\VS{11}Sua cabeça é
{\ADD{como}} o mais fino ouro; seus cabelos são crespos, pretos como o corvo.
\VS{12}Seus olhos são como pombas junto às correntes de águas, lavados em leite, colocados como
{\ADD{joias}} .
\VS{13}As laterais de seu rosto como um canteiro de especiarias,
{\ADD{como}} caixas aromáticas; seus lábios são
{\ADD{como}} lírios, que gotejam fragrante mirra.
\VS{14}Suas mãos são
{\ADD{como}} anéis de ouro com pedras de crisólitos; seu abdome
{\ADD{como}} o brilhante marfim, coberto de safiras.
\VS{15}Suas pernas são
{\ADD{como}} colunas de mármore, fundadas sobre bases de ouro puro; seu aspecto é como o Líbano, precioso como os cedros.
\VS{16}Sua boca é
{\ADD{cheia}} de doçura, e ele é desejável em todos os sentidos; assim é o meu amado; assim é o meu querido, ó filhas de Jerusalém.

\par }\Chap{6}{\PP \VerseOne{1}{\ADD{Moças}} : Para onde foi o teu amado, ó tu mais bela entre as mulheres? Para que direção se virou o teu amado, para o procurarmos contigo?
\VS{2}{\ADD{Ela}} : Meu amado desceu ao seu jardim, aos canteiros de especiarias, para apascentar
{\ADD{seu rebanho}} nos jardins, e para colher lírios.
\VS{3}Eu sou do meu amado, e meu amado é meu; ele apascenta entre os lírios.
\VS{4}{\ADD{Ele}} : Tu és bela, minha querida, como Tirza, agradável como Jerusalém; és formidável como bandeiras
{\ADD{de exércitos}} .
\VS{5}Afasta teus olhos de mim, pois eles me deixam desconcertado. Teu cabelo é como um rebanho de cabras, que descem de Gileade.
\VS{6}Teus dentes são como um rebanho de ovelhas, que sobem do lavatório; todas produzem gêmeos, e não há estéril entre elas.
\VS{7}Como um pedaço de romã, assim são as laterais de teu rosto abaixo de teu véu.
\VS{8}Sessenta são as rainhas, e oitenta as concubinas; e as donzelas são inúmeras;
\VS{9}{\ADD{Porém}} uma é a minha pomba, minha perfeita, a única de sua mãe, a mais querida daquela que a gerou. As moças
\FTNT{*}{{\FR 6:9 }
moças = lit. filhas
} a viram, e a chamaram de bem-aventurada; as rainhas e as concubinas a elogiaram.
\VS{10}Quem é esta, que aparece como o nascer do dia, bela como a lua, brilhante como o sol, formidável como bandeiras
{\ADD{de exércitos}} ?
\VS{11}Desci ao jardim das nogueiras, para ver os frutos do vale; para ver se as videiras estavam floridas,
{\ADD{e se}} as romãzeiras brotavam.
\VS{12}Sem eu perceber, minha alma me pôs nas carruagens de meu nobre povo.
\VS{13}{\ADD{Moças}} : Volta! Volta, Sulamita! Volta! Volta, e nós te veremos!
{\ADD{Ele}} : Por que
{\ADD{quereis}} ver a Sulamita, como a dança
\FTNT{†}{{\FR 6:13 }
dança
trad. alt. fileira
} de duas companhias?
\FTNT{‡}{{\FR 6:13 }
duas companhias
trad. alt. dois exércitos - ou transliteração Maanaim
}

\par }\Chap{7}{\PP \VerseOne{1}{\ADD{Ele}} : Como são belos os teus pés nas sandálias, ó filha de príncipe! Os contornos de tuas coxas são como joias,
{\ADD{como}} obra das mãos de um artesão.
\VS{2}Teu umbigo é
{\ADD{como}} uma taça redonda, que não falta bebida; teu abdome é
{\ADD{como}} um amontoado de trigo, rodeado de lírios.
\VS{3}Teus dois seios são como dois filhos gêmeos da corça.
\VS{4}Teu pescoço é como uma torre de marfim; teus olhos são
{\ADD{como}} os tanques de peixes de Hesbom, junto à porta de Bate-Rabim. Teu nariz é como a torre do Líbano, que está de frente a Damasco.
\VS{5}Tua cabeça sobre ti, como o
{\ADD{monte}} Carmelo, e o trançado dos cabelos cabeça como púrpura; o rei está
{\ADD{como que}} atado em
{\ADD{tuas}} tranças.
\VS{6}Como tu és bela! Como tu és agradável, ó amor em delícias!
\VS{7}Esta tua estatura é semelhante à palmeira, e teus seios são
{\ADD{como}} cachos
{\ADD{de uvas}} .
\VS{8}Eu disse: Subirei à palmeira, pegarei dos seus ramos; e então teus seios serão como os cachos da videira, e a fragrância de teu nariz como o das maçãs.
\VS{9}E tua boca
{\ADD{seja}} como o bom vinho;
{\ADD{Ela}} :
{\ADD{Vinho}} que se entra a meu amado suavemente, e faz falarem os lábios dos que dormem.
\VS{10}Eu sou do meu amado, e ele me deseja.
\VS{11}Vem, amado meu! Saiamos ao campo, passemos as noites nas aldeias.
\VS{12}Saiamos de madrugada até às vinhas, vejamos se as videiras florescem, se suas flores estão se abrindo, se as romãzeiras já estão brotando; ali te darei os meus amores.
\VS{13}As mandrágoras espalham seu perfume, e junto a nossas portas há todo tipo de excelentes frutos, novos e velhos; eu os guardei para ti, meu amado.

\par }\Chap{8}{\PP \VerseOne{1}Como gostaria que tu fosses como meu irmão, que mamava os seios de minha mãe. Quando eu te achasse na rua, eu te beijaria, e não me desprezariam.
\VS{2}Eu te levaria e te traria à casa de minha mãe, aquela que me ensinou, e te daria a beber do vinho aromático, do suco de minhas romãs.
\VS{3}Esteja sua mão esquerda debaixo de minha cabeça, e sua direita me abrace.
\VS{4}Eu vos ordeno, filhas de Jerusalém: jurai que não acordeis, nem desperteis ao amor, até que ele queira.
\VS{5}{\ADD{Outros}} : Quem é esta, que sobe da terra desabitada, reclinada sobre seu amado?
{\ADD{Ela}} : Debaixo de uma macieira te despertei; ali tua mãe te teve com dores; ali te deu à luz aquela que te gerou.
\VS{6}Põe-me como selo sobre o teu coração, como selo sobre o teu braço; porque o amor é forte como a morte, e o ciúme
\FTNT{*}{{\FR 8:6 }
Ou: paixão
} é severo como o Xeol; suas brasas são brasas de fogo, labaredas intensas.
\FTNT{†}{{\FR 8:6 }
labaredas intensas
Trad. alternativa: labaredas do SENHOR
}
\VS{7}Muitas águas não poderiam apagar este amor, nem os rios afogá-lo; ainda que alguém desse todos os bens de sua casa por este amor, certamente o desprezariam.
\VS{8}{\ADD{Irmãos dela}} : Temos uma irmã pequena, que ainda não têm seios; que faremos a esta nossa irmã, no dia que dela falarem?
\VS{9}Se ela for um muro, edificaremos sobre ela um palácio de prata; e se ela for uma porta, então a cercaremos com tábuas de cedro.
\VS{10}{\ADD{Ela}} : Eu sou um muro, e meus seios como torres; então eu era aos olhos dele, como aquela que encontra a paz.
\VS{11}Salomão teve uma vinha em Baal-Hamom; ele entregou esta vinha a uns guardas, e cada um
{\ADD{lhe}} trazia como
{\ADD{faturamento}} de seus frutos mil
{\ADD{moedas}} de prata.
\VS{12}Minha vinha, que me pertence, está à minha disposição;
\FTNT{‡}{{\FR 8:12 }
à minha disposição = lit. perante minha face
} as mil
{\ADD{moedas}} de prata são para ti, Salomão, e duzentas para os guardas de teu fruto.
\VS{13}{\ADD{Ele}} : Ó tu que habitas nos jardins, teus colegas prestam atenção à tua voz: faze com que eu
{\ADD{também possa}} ouvi-la.
\VS{14}{\ADD{Ela}} : Vem depressa, meu amado! E faze-te semelhante ao corço, ao filhote de cervos, nas montanhas de ervas aromáticas!\par }